\section{Local regularization on normal varieties}\label{sec:local-regularization}

\begin{convention}\label{conv:sec5}
Throughout, $k$ is an algebraically closed field with the trivial absolute value. From
\S\ref{subsec:polyconvex} on, $k$ is also countable. An \emph{affine variety} is $Y=\Spec A$
with $A$ a finitely generated $k$-algebra that is a domain; it is \emph{normal} if $A$ is
integrally closed. For $f\in A$ we write $Y_f=\Spec A_f$.
\end{convention}

\subsection{The analytification of an affine variety}\label{subsec:berk}

\begin{definition}[Affine analytification]\label{def:berk}
\lean{EnvelopeProperty.Berk, EnvelopeProperty.Berk.logAbs}
\leanok
The analytification $Y^\an$ of $Y=\Spec A$ is the set of maps $x\colon A\to\R_{\ge0}$,
$f\mapsto|f|_x$, with $|0|_x=0$, $|1|_x=1$, $|fg|_x=|f|_x|g|_x$,
$|f+g|_x\le\max\{|f|_x,|g|_x\}$ and $|c|_x=1$ for $c\in k^\times$, with the coarsest
topology making all $x\mapsto|f|_x$ continuous. We put $\log|f|(x)=\log|f|_x\in[-\infty,\infty)$.
In the additive notation of \cite{FHX}, $x$ is the semivaluation $f\mapsto-\log|f|_x$.
\end{definition}

\begin{definition}[Trivial valuation and scaling {\cite[\S2.1]{FHX}}]\label{def:triv-scaling}
\lean{EnvelopeProperty.Berk.triv, EnvelopeProperty.Berk.rpow, EnvelopeProperty.Berk.zeroScale}
\leanok
\uses{def:berk}
The trivial valuation $\vtriv$ has $|f|=1$ for $f\ne0$. For $t>0$ the point $t\cdot x$ has
$|f|_{t\cdot x}=|f|_x^t$, and $0\cdot x$ has $|f|_{0\cdot x}=1$ if $|f|_x\ne0$ and $0$ otherwise.
\end{definition}

\begin{lemma}[Compact polydiscs]\label{lem:polydisc-compact}
\lean{EnvelopeProperty.Berk.abs_le_max_pow_of_adjoin_eq_top, EnvelopeProperty.Berk.isClosedEmbedding_toFun, EnvelopeProperty.Berk.isCompact_setOf_forall_abs_le_finset, EnvelopeProperty.Berk.instLocallyCompactSpace, EnvelopeProperty.Berk.instSigmaCompactSpace}
\leanok
\uses{def:berk}
Let $z_1,\dots,z_n$ generate $A$. For every $f\in A$ there is $d\in\N$ with
$|f|_x\le\max\{1,|z_1|_x,\dots,|z_n|_x\}^d$ for all $x$. For $R\ge0$ the polydisc
$\{x:|z_l|_x\le R\ \forall l\}$ is compact. Consequently $Y^\an$ is Hausdorff, locally compact and
$\sigma$-compact, and second countable when $k$ is countable.
\end{lemma}

\begin{proof}
Write $f=\sum_\alpha c_\alpha z^\alpha$ with $|\alpha|\le d$; since $|c_\alpha|_x\le1$, the
ultrametric inequality gives the bound. The map $x\mapsto(|f|_x)_{f\in A}$ is a homeomorphism
onto the subset of $\R_{\ge0}^A$ cut out by the closed conditions of Definition~\ref{def:berk},
which is closed; the polydisc maps into the product of the compact intervals
$[0,\max\{1,R\}^{d_f}]$, so it is compact by Tychonoff's theorem. Every point $x$ has the
neighbourhood $\{|z_l|<|z_l|_x+1\}$ inside a compact polydisc, $Y^\an$ is the union of the
polydiscs of radius $m\in\N$, and if $k$ (hence $A$) is countable then $\R_{\ge0}^A$ is second
countable.
\end{proof}

\begin{lemma}[Principal opens {\cite[\S5]{FHX}}]\label{lem:principal-open}
\lean{EnvelopeProperty.Berk.isOpenEmbedding_awayComap, EnvelopeProperty.Berk.range_awayComap, EnvelopeProperty.Berk.awayComap_val_mk'}
\leanok
\uses{def:berk,spectrum:localization-trivial-field}
For $f\in A$ the restriction map $Y_f^\an\to Y^\an$ is an open embedding with image
$\{x:|f|_x\ne0\}$, and $|a/f^n|=|a|/|f|^n$ on it.
\end{lemma}

\begin{proof}
A point $x$ with $|f|_x\ne0$ extends uniquely to $A_f$ by $|a/f^n|=|a|_x/|f|_x^n$; this formula
is continuous on the open set $\{|f|\ne0\}$, which gives continuity of the inverse.
\end{proof}

\subsection{Local psh functions}\label{subsec:local-psh}

\begin{definition}[Convex PL and locally psh functions {\cite[Definition~5.1]{FHX}}]\label{def:convex-pl}
\lean{EnvelopeProperty.IsConvexPLOn, EnvelopeProperty.IsDecrConvexPLLimitOn, EnvelopeProperty.IsLocallyPSHOn}
\leanok
\uses{def:berk}
Let $U\subseteq Y^\an$ be open. A function $u$ is \emph{convex PL} on $U$ if every point of $U$
has an open neighbourhood $W\subseteq U$ with elements $f_1,\dots,f_r\in A$ not vanishing on $W$
and finitely many (at least one) $(a_j,b_j)\in\Q^r\times\Q$ with
$u=\max_j(\sum_ia_{ij}\log|f_i|+b_j)$ on $W$. It is \emph{locally psh} if every point has an open
neighbourhood on which it is the limit of a decreasing sequence of convex PL functions.
\end{definition}

\begin{remark}\label{rmk:convex-pl-affine}
\uses{def:convex-pl}
In \cite[Definition~5.1]{FHX} the $f_i$ are invertible functions on a Zariski open
$V_0\supseteq W$ and $V$ may be any scheme of finite type. For $V=Y$ affine this is equivalent:
near any point of $V_0$ an invertible function is a quotient $g/h^n$ of elements of $A$ not
vanishing there, and conversely elements of $A$ not vanishing on $W$ are invertible on the
principal open defined by their product. The general schemes and base fields of
\cite[Definition~5.1]{FHX} are needed only in \S\ref{subsec:norms} and Corollary~\ref{cor:normal-weak-global}.
\end{remark}

\begin{lemma}[Operations on convex PL functions {\cite[Lemma~5.2]{FHX}}]\label{lem:local-basic}
\lean{EnvelopeProperty.IsConvexPLOn.max, EnvelopeProperty.IsConvexPLOn.add, EnvelopeProperty.IsConvexPLOn.smul, EnvelopeProperty.IsConvexPLOn.add_const, EnvelopeProperty.IsConvexPLOn.add_smul_logAbs, EnvelopeProperty.IsConvexPLOn.piecewise_max}
\leanok
\uses{def:convex-pl}
Finite maxima and sums, positive rational multiples and rational translates of convex PL
functions are convex PL, and so is $u+q\log|f|$ on $U\cap\{|f|\ne0\}$ for $q\in\Q$. If
$U'\subseteq U$ are open, $u$ is convex PL on $U$, $v$ is convex PL on $U'$ and the closure in $U$
of $\{x\in U':v(x)>u(x)\}$ lies in $U'$, then the function equal to $\max\{u,v\}$ on $U'$ and to $u$
on $U\setminus U'$ is convex PL on $U$.
\end{lemma}

\begin{proof}
A sum of two finite maxima is the maximum of the pairwise sums; the other operations act on the
coefficients. For the gluing, on $U'$ use the maximum; a point of $U\setminus U'$ has a
neighbourhood disjoint from the closure in question, on which the function equals $u$.
\end{proof}

\subsection{Regularizable functions on affine varieties}\label{subsec:regularizable}

\begin{definition}[Affine Fubini--Study and regularizable functions {\cite[Definition~5.7]{FHX}}]\label{def:fs}
\lean{EnvelopeProperty.FSData, EnvelopeProperty.IsFS, EnvelopeProperty.HasLogGrowth, EnvelopeProperty.IsRegularizable}
\leanok
\uses{def:berk}
An \emph{affine Fubini--Study function} is
$g=\max\{a_1\log|f_1|+b_1,\dots,a_r\log|f_r|+b_r,b_0\}$ with $a_i\in\Q_{>0}$, $b_i\in\Q$ and
$f_1,\dots,f_r\in A$ including a set of generators of the $k$-algebra $A$; $\FS(Y)$ is the set of
these. A function $u$ has \emph{logarithmic growth} if $u\le g+C$ for some $g\in\FS(Y)$ and
$C\in\R$. It is \emph{regularizable}, $u\in\Rcal(Y)$, if $u\not\equiv-\infty$ and $u$ is the
limit of a decreasing sequence in $\FS(Y)$.
\end{definition}

\begin{lemma}[Elementary properties of FS functions {\cite[Definition~5.7]{FHX}}]\label{lem:fs-basic}
\lean{EnvelopeProperty.FSData.continuous_eval, EnvelopeProperty.FSData.const_le_eval, EnvelopeProperty.FSData.isConvexPLOn_univ, EnvelopeProperty.FSData.eval_add, EnvelopeProperty.FSData.eval_sup, EnvelopeProperty.FSData.eval_smul, EnvelopeProperty.FSData.eval_addConst, EnvelopeProperty.FSData.eval_mapSup}
\leanok
\uses{def:fs, def:convex-pl, lem:local-basic}
Functions in $\FS(Y)$ are continuous, bounded below and convex PL on $Y^\an$. $\FS(Y)$ is
stable under finite maxima, sums, multiplication by positive rationals and addition of rational
constants. If $\pi\colon Y'\to Y$ is a morphism, $g\in\FS(Y)$ and $g'\in\FS(Y')$, then
$\max\{g\circ\pi,g'\}\in\FS(Y')$.
\end{lemma}

\begin{proof}
$g\ge b_0$. Near a point $x$, a term with $|f_i|_x=0$ equals $-\infty$ at $x$, hence is $<b_0$
nearby, and may be omitted; the remaining $f_i$ do not vanish near $x$. Sums: a sum of maxima
is the maximum of the pairwise sums, and $a\log|f|+a'\log|f'|=N^{-1}\log|f^{Na}f'^{Na'}|$ for a
common denominator $N$. For $\max\{g\circ\pi,g'\}$ the generators come from $g'$.
\end{proof}

\begin{lemma}[Growth estimates {\cite[Lemma~5.8]{FHX}}]\label{lem:growth}
\lean{EnvelopeProperty.FSData.exists_eval_comap_le, EnvelopeProperty.FSData.exists_eval_le_comap_of_isIntegral, EnvelopeProperty.valuation_le_of_monic_eval₂_eq_zero}
\leanok
\uses{def:fs, lem:polydisc-compact}
(1) If $\pi\colon Y'\to Y$ is a morphism of affine varieties, $g\in\FS(Y)$ and $g'\in\FS(Y')$,
there are $C\ge0$, $C'$ with $g\circ\pi\le Cg'+C'$. (2) If $\pi\colon Y\to S$ is finite and
surjective, $g\in\FS(Y)$ and $h\in\FS(S)$, there are $C,C'$ with $g\le C\,h\circ\pi+C'$.
\end{lemma}

\begin{proof}
(1) By Lemma~\ref{lem:polydisc-compact}, $\log|f|\le d\max\{0,\log|z_1|,\dots,\log|z_N|\}$ for
generators $z_l$ of $A'$ occurring in $g'$; apply this to the pullbacks of the functions defining
$g$, and bound $\max\{0,\log|z_l|\}$ by a multiple of $g'$ plus a constant. (2) Each $f\in A$
satisfies $f^m+b_{m-1}f^{m-1}+\dots+b_0=0$ with $b_i$ from the coordinate ring of $S$. If
$|f|>|b_i|^{1/(m-i)}$ for all $i$ then $|f^m|>|b_if^i|$ for all $i$, contradicting the
ultrametric inequality; hence $\log|f|\le\max_i(m-i)^{-1}\log|b_i|\circ\pi$. Apply (1) on $S$.
\end{proof}

\begin{theorem}[Growth-controlled approximation; formerly imported {\cite[Lemma~5.11]{FHX}, \cite[Lemma~3.25]{Wu26}}]\label{thm:import-growth-control}
\lean{EnvelopeProperty.WuRegularization, EnvelopeProperty.ConvexPLApproximation.growthControl, EnvelopeProperty.LocalPositivity.toConvexPLApproximation}
\leanok
\uses{def:fs, thm:wu-growth-control, thm:wu-pl-approx, lem:wu-local-positivity}
Let $Y$ be normal, $u\in\Rcal(Y)$ and $g\in\FS(Y)$ with $u\le g+C$. For every rational $c>1$
there are $u_j\in\FS(Y)$ decreasing to $u$ with $u_j-cg$ bounded above.
\end{theorem}

\begin{theorem}[Local-to-global regularization; formerly imported {\cite[Lemma~5.11]{FHX}, \cite[Lemma~3.26]{Wu26}}]\label{thm:import-local-to-global}
\lean{EnvelopeProperty.WuRegularization, EnvelopeProperty.ConvexPLApproximation.localToGlobal, EnvelopeProperty.ConvexPLApproximation.toWuRegularization}
\leanok
\uses{def:fs, def:convex-pl, thm:wu-local-to-global, thm:wu-pl-approx, lem:wu-local-positivity}
Let $Y$ be normal and $v\not\equiv-\infty$ of logarithmic growth such that every compact subset
of $Y^\an$ has an open neighbourhood on which $v$ is the limit of a decreasing sequence of convex
PL functions. Then $v\in\Rcal(Y)$.
\end{theorem}

\begin{remark}\label{rmk:imports}
\uses{thm:import-growth-control, thm:import-local-to-global}
These are \cite[Lemma~5.11(5),(6)]{FHX}, i.e.\ \cite[Lemmas~3.25(2), 3.26]{Wu26} with the
corrections of \cite{FHX}; statement (6) is stronger than Wu's literal Lemma~3.26 (arbitrary local
convex PL charts), the extra strength coming from \cite[Lemma~5.5]{FHX}. They enter the formalization
as an explicit hypothesis and are not formalized, so every node using them is proved only
conditionally on these two statements. Their proofs pass to a normal projective compactification
\cite[Lemmas~5.9, 5.10]{FHX} and use Boucksom--Jonsson's global theory \cite{BJ22} (PSH2,
countable regularization in an ample class, test configurations) together with the positivity
criterion of Chambert-Loir--Ducros \cite{CLD12} \cite[Lemma~5.5]{FHX}. Only the upper bound in
Theorem~\ref{thm:import-growth-control} is used.
\end{remark}

\begin{lemma}[Finiteness at the trivial valuation {\cite[Lemma~5.11]{FHX}}]\label{lem:reg-triv-finite}
\lean{EnvelopeProperty.IsRegularizable.triv_ne_bot, EnvelopeProperty.FSData.two_mul_eval_le, EnvelopeProperty.FSData.eval_zeroScale_le_triv}
\leanok
\uses{def:fs, def:triv-scaling}
Every $u\in\Rcal(Y)$ is finite at $\vtriv$.
\end{lemma}

\begin{proof}
Let $u_j\in\FS(Y)$ decrease to $u$. For $x\in Y^\an$ the function
$t\mapsto u_j(t\cdot x)=\max\{b_0,\ a_it\log|f_i|_x+b_i\}$ on $[0,\infty)$ is a finite maximum of
affine functions (terms with $|f_i|_x=0$ being omitted), hence convex, so
$u_j(x)\le\frac12(u_j(0\cdot x)+u_j(2\cdot x))$. Moreover $u_j(0\cdot x)\le u_j(\vtriv)$, since
$|f|_x\ne0$ forces $f\ne0$. Letting $j\to\infty$,
$u(x)\le\frac12(u(\vtriv)+u(2\cdot x))$ with $u(2\cdot x)<\infty$. If $u(\vtriv)=-\infty$ then
$u\equiv-\infty$, which is excluded. Finally $u(\vtriv)\le u_1(\vtriv)<\infty$.
\end{proof}

\begin{remark}\label{rmk:triv-finite-scope}
\uses{lem:reg-triv-finite}
\cite[Lemma~5.11(1)]{FHX} asserts finiteness on all of $Y^\divv$; Section~5 only uses
finiteness at $\vtriv$, which is proved here directly.
\end{remark}

\begin{lemma}[Elementary closure properties {\cite[Lemma~5.11]{FHX}}]\label{lem:reg-closure-elementary}
\lean{EnvelopeProperty.isRegularizable_finsetSup_fsTerm, EnvelopeProperty.IsRegularizable.sup, EnvelopeProperty.IsRegularizable.add, EnvelopeProperty.IsRegularizable.smul, EnvelopeProperty.IsRegularizable.add_const, EnvelopeProperty.IsRegularizable.comp_comap, EnvelopeProperty.IsRegularizable.comp_awayComap}
\leanok
\uses{lem:fs-basic, lem:reg-triv-finite}
(3) If $q=\max_i(a_i\log|f_i|+b_i)\not\equiv-\infty$ with $a_i\in\Q_{>0}$, $b_i\in\Q$, then
$q\in\Rcal(Y)$; $\Rcal(Y)$ is stable under finite maxima, finite sums, positive rational multiples
and real translates. (4) For a morphism $\pi\colon Y'\to Y$ and $u\in\Rcal(Y)$, $u\circ\pi$ is
identically $-\infty$ or regularizable.
\end{lemma}

\begin{proof}
(3) With $b=\max\{0,\log|z_1|,\dots,\log|z_n|\}$ for generators $z_l$, the functions
$\max\{q,b/j-j\}\in\FS(Y)$ decrease to $q$ since $b\ge0$ is finite. Maxima and sums of
approximants approximate maxima and sums, and the limits are not $\equiv-\infty$ by
Lemma~\ref{lem:reg-triv-finite}; for $u+c$ use $u_j+c_j$ with rational $c_j\downarrow c$.
(4) If $u_j\downarrow u$ then $\max\{u_j\circ\pi,b'/j-j\}\in\FS(Y')$ decrease to $u\circ\pi$.
\end{proof}

\begin{lemma}[Diagonal approximation on compact sets]\label{lem:diagonal-compact}
\lean{EnvelopeProperty.exists_antitone_fs_of_antitone_regularizable, EnvelopeProperty.exists_antitone_diagonal}
\leanok
\uses{lem:fs-basic}
Let $K\subseteq Y^\an$ be compact and $u_i\in\Rcal(Y)$ decrease to $u$. There are $h_m\in\FS(Y)$,
decreasing on $K$, with $\lim_mh_m=u$ on $K$.
\end{lemma}

\begin{proof}
Let $u_i=\lim_j\downarrow g_{ij}$ with $g_{ij}\in\FS(Y)$, and enumerate the countable family
$T=\{g_{ij}+q:i,j\in\N,\ q\in\Q_{>0}\}$ as $t_1,t_2,\dots$; then $\inf T=u$. Given
$s=g_{i_1j_1}+q_1$ and $t=g_{i_2j_2}+q_2$ in $T$, choose $i\ge i_1,i_2$ and rational
$0<q<\min\{q_1,q_2\}$. Since the $g$'s are finite, $u_i+q<\min\{s,t\}$ everywhere, so the open
sets $\{g_{ij}+q<\min\{s,t\}\}$ increase in $j$ and cover $K$; by compactness one of them contains
$K$. Starting from $h_1=t_1$, this gives inductively $h_{m+1}\in T$ with
$h_{m+1}\le\min\{h_m,t_{m+1}\}$ on $K$. Then $u\le h_m\le t_m$ on $K$ and $\inf_mh_m=u$ there.
\end{proof}

\begin{lemma}[Decreasing limits {\cite[Lemma~5.11]{FHX}}]\label{lem:reg-decreasing}
\lean{EnvelopeProperty.IsRegularizable.of_antitone}
\leanok
\uses{lem:diagonal-compact, thm:import-local-to-global, lem:polydisc-compact}
If $Y$ is normal and $u_i\in\Rcal(Y)$ decrease to $u\not\equiv-\infty$, then $u\in\Rcal(Y)$.
\end{lemma}

\begin{proof}
Every compact $K$ has a compact neighbourhood $K'$ (Lemma~\ref{lem:polydisc-compact});
Lemma~\ref{lem:diagonal-compact} on $K'$ exhibits $u$ on $\interior K'$ as the limit of a
decreasing sequence of FS functions, which are convex PL. Since $u\le u_1$ has logarithmic growth,
Theorem~\ref{thm:import-local-to-global} applies.
\end{proof}

\begin{lemma}[Principal-open extension {\cite[Lemma~5.12]{FHX}}]\label{lem:principal-extension}
\lean{EnvelopeProperty.IsRegularizable.awayExtend_add, EnvelopeProperty.IsConvexPLOn.awayExtend, EnvelopeProperty.isConvexPLOn_univ_max_awayExtend, EnvelopeProperty.FSData.eval_awayEps}
\leanok
\uses{lem:principal-open, thm:import-growth-control, thm:import-local-to-global, lem:reg-decreasing, lem:reg-triv-finite, lem:local-basic, lem:growth}
Let $0\ne f\in A$ and $u\in\Rcal(Y_f)$ with $u\le g+C$ on $Y_f^\an$ for some $g\in\FS(Y)$. For
every $c\in\Q_{>0}$ the function equal to $u+c\log|f|$ on $Y_f^\an$ and to $-\infty$ on
$\{|f|=0\}$ belongs to $\Rcal(Y)$.
\end{lemma}

\begin{proof}
Assume $g\ge0$. Choose a rational $0<\epsilon<c/2$; $g_\epsilon=\max\{g,-\epsilon\log|f|\}$ is in
$\FS(Y_f)$ and $u\le g_\epsilon+C$. By Theorem~\ref{thm:import-growth-control} there are
$u_j\in\FS(Y_f)$ decreasing to $u$ with $u_j\le2g_\epsilon+C_j$. Let $v_j=\max\{u_j+c\log|f|,-j\}$ on
$Y_f^\an$ and $v_j=-j$ where $|f|=0$. For fixed $j$,
$u_j+c\log|f|\le\max\{2g+c\log|f|,(c-2\epsilon)\log|f|\}+C_j$ tends to $-\infty$ at each point of
$\{|f|=0\}$, so $v_j=-j$ near such points, and $v_j$ is convex PL on $Y^\an$ (on $Y_f^\an$ the
functions $\log|a/f^n|=\log|a|-n\log|f|$ are of the required form). The same bound and
Lemma~\ref{lem:growth} give logarithmic growth on $Y$, so $v_j\in\Rcal(Y)$ by
Theorem~\ref{thm:import-local-to-global}. The $v_j$ decrease to the asserted extension, which is
finite at $\vtriv$ by Lemma~\ref{lem:reg-triv-finite}; apply Lemma~\ref{lem:reg-decreasing}.
\end{proof}

\subsection{Polynomial convexity}\label{subsec:polyconvex}

\begin{definition}[Polynomial hulls {\cite[\S5.3]{FHX}}]\label{def:hull}
\lean{EnvelopeProperty.hull, EnvelopeProperty.IsPolyConvex, EnvelopeProperty.IsRunge, EnvelopeProperty.SatisfiesHullIneq}
\leanok
\uses{def:berk}
For compact $K\subseteq Y^\an$, $\widehat K=\{y:\log|f|(y)\le\sup_K\log|f|\ \forall f\in A\}$. $K$ is
polynomially convex if $\widehat K=K$; an open $D$ is Runge if $\widehat K\subseteq D$ for all compact
$K\subseteq D$; $u$ satisfies the hull inequality if $u(y)\le\sup_Ku$ for all compact $K$ and
$y\in\widehat K$.
\end{definition}

\begin{lemma}[Hulls and the hull inequality {\cite[\S5.3]{FHX}}]\label{lem:hull-basic}
\lean{EnvelopeProperty.isCompact_hull, EnvelopeProperty.subset_hull, EnvelopeProperty.hull_hull, EnvelopeProperty.FSData.satisfiesHullIneq, EnvelopeProperty.iInf_iSup_eq_iSup_iInf_of_antitone, EnvelopeProperty.IsRegularizable.satisfiesHullIneq, EnvelopeProperty.isRunge_iUnion_of_monotone}
\leanok
\uses{def:hull, def:fs, lem:polydisc-compact}
For compact $K$, $\widehat K$ is a compact set containing $K$ and $\widehat{\widehat K}=\widehat K$. Every
function in $\FS(Y)$ and every $u\in\Rcal(Y)$ satisfies the hull inequality. An increasing union of
Runge open sets is Runge.
\end{lemma}

\begin{proof}
$\widehat K$ is closed and lies in the polydisc $\{|z_l|\le\sup_K|z_l|\}$. Each term
$a\log|f|+b$ of an FS function satisfies the inequality, hence so does their maximum. If
$u_j\in\FS(Y)$ decrease to $u$ then $u(y)\le u_j(y)\le\sup_Ku_j$, and $\sup_Ku_j\downarrow\sup_Ku$:
for $t>\sup_Ku$ the open sets $\{u_j<t\}$ increase and cover $K$. A compact subset of an
increasing union of open sets lies in one of them.
\end{proof}

\begin{lemma}[Polyhedral neighbourhoods {\cite[Lemma~5.14]{FHX}}]\label{lem:polyhedral-neighborhoods}
\lean{EnvelopeProperty.exists_fs_polyhedral_nhd, EnvelopeProperty.IsRunge.exists_fs_exhaustion, EnvelopeProperty.exists_lt_logAbs_of_notMem_hull}
\leanok
\uses{lem:hull-basic, lem:fs-basic}
(1) If $H$ is a polynomially convex compact set and $A'\supseteq H$ open, there is $g\in\FS(Y)$ with
$g\ge0$, $H\subseteq\interior\{g=0\}$ and $\{g\le1\}\subseteq A'$. (2) Every Runge open $D$ is the
union of an increasing sequence $\{g_j<1\}$ with $g_j\in\FS(Y)$.
\end{lemma}

\begin{proof}
(1) Choose generators $z_l$ and $R\in\Q$ with $H\subseteq\{\max_l\log|z_l|<R\}$. The set
$E=\{\max_l\log|z_l|\le R+1\}\setminus A'$ is compact and disjoint from $\widehat H=H$; for $y\in E$,
replacing a function by a power, there are $f_y$ and $a_y\in\Q$ with
$\sup_H\log|f_y|<a_y<\log|f_y(y)|-1$. Finitely many $\{\log|f_{y_\ell}|>a_{y_\ell}+1\}$ cover $E$.
Put $g=\max\{0,\max_l(\log|z_l|-R),\max_\ell(\log|f_{y_\ell}|-a_{y_\ell})\}$. Then $g=0$ on the open
set where all non-constant terms are negative, which contains $H$, and $g>1$ off $A'$.
(2) Write $D=\bigcup_jK_j$ with $K_j$ compact increasing (Lemma~\ref{lem:polydisc-compact},
$k$ countable). With $Q_0=\varnothing$, apply (1) to $H_j=\widehat{K_j\cup Q_{j-1}}\subseteq D$ and
$A'=D$, and put $Q_j=\{g_j\le1\}$, which is compact. Then
$K_j\cup Q_{j-1}\subseteq H_j\subseteq\{g_j<1\}\subseteq Q_j\subseteq D$.
\end{proof}

\begin{lemma}[Rossi's local maximum principle]\label{lem:rossi}
\lean{EnvelopeProperty.mem_hull_frontier_inter_hull}
\leanok
\uses{lem:polyhedral-neighborhoods, lem:local-basic, thm:import-local-to-global, lem:hull-basic}
Let $Y$ be normal, $K$ compact, $x\in\widehat K$ and $V$ a relatively compact open
neighbourhood of $x$ with $\overline V\cap K=\varnothing$. Then $x$ lies in the hull of
$\partial V\cap\widehat K$. Here $x\notin K$ follows from $x\in V$ and $\overline V\cap K=\varnothing$.
\end{lemma}

\begin{proof}
Otherwise there are $f\in A$ and $t\in\Q$ with $\sup_{\partial V\cap\widehat K}\log|f|<t<\log|f(x)|$.
The compact set $E=\partial V\cap\{\log|f|\ge t\}$ is disjoint from $\widehat K$; by
Lemma~\ref{lem:polyhedral-neighborhoods}(1) choose $g\in\FS(Y)$, $g\ge0$, $g=0$ near $\widehat K$ and
$g>1$ on $E$. Choose an integer $C>0$ with $C>\sup_{\overline V}(\log|f|-t)$ and let
$v=\max\{Cg,\log|f|-t\}$ on $V$ and $v=Cg$ outside $V$. On $\partial V$, $\log|f|-t<Cg$, so $v$ is
convex PL by Lemma~\ref{lem:local-basic}; it is bounded by an FS function, so
$v\in\Rcal(Y)$ by Theorem~\ref{thm:import-local-to-global}. But $v=0$ on $K$ and $v(x)>0$,
contradicting Lemma~\ref{lem:hull-basic}.
\end{proof}

\begin{lemma}[Gauss points]\label{lem:gauss-point}
\lean{EnvelopeProperty.Berk.gauss, EnvelopeProperty.Berk.proj_gauss, EnvelopeProperty.Berk.gauss_val_X, EnvelopeProperty.Berk.val_le_gauss_val, EnvelopeProperty.Berk.continuous_gauss, EnvelopeProperty.Berk.logAbs_gauss_lt_iff}
\leanok
\uses{def:berk,spectrum:gauss-valuation,spectrum:gauss-joint-continuity,spectrum:trivial-field-bridge}
Let $\widetilde Y=Y\times\A^1=\Spec A[w]$. For $y\in Y^\an$ and $r\ge0$ the formula
$|\sum_\ell a_\ell w^\ell|=\max_\ell|a_\ell|_yr^\ell$ defines a point $\eta_{y,r}\in\widetilde Y^\an$
over $y$ with $|w|=r$; it depends continuously on $y$, and every $z$ over $y$ with $|w|_z\le r$
satisfies $|F|_z\le|F|_{\eta_{y,r}}$ for all $F\in A[w]$.
\end{lemma}

\begin{proof}
The general Gauss seminorm construction proves multiplicativity directly over a
commutative ring. Specializing to the trivial-field relative spectrum gives the
stated point, maximality, and continuity through the evaluation-preserving comparison.
\end{proof}

\begin{lemma}[Hartogs radius {\cite[Lemma~5.16]{FHX}}]\label{lem:hartogs}
\lean{EnvelopeProperty.IsRegularizable.of_isRunge_hartogs, EnvelopeProperty.Hartogs.gauss_rad_mem_iff, EnvelopeProperty.Hartogs.isConvexPLOn_radiusFun}
\leanok
\uses{lem:gauss-point, spectrum:gauss-compact-fibre-disc, lem:polyhedral-neighborhoods, lem:local-basic, thm:import-local-to-global}
Let $Y$ be normal and $u$ upper semicontinuous, of logarithmic growth and not identically
$-\infty$. If $D=\{u\circ p+\log|w|<0\}\subseteq\widetilde Y^\an$ is Runge, then $u\in\Rcal(Y)$.
\end{lemma}

\begin{proof}
$D$ is open. By Lemma~\ref{lem:polyhedral-neighborhoods}(2) on $\widetilde Y$, $D$ is the
increasing union of $W_j=\{g_j<1\}\cap\{\log|w|<j\}$, and each $\{g_j<1\}$ (if nonempty) is
$\bigcap_\alpha\{\log|F_{j\alpha}|<c_{j\alpha}\}$ with $F_{j\alpha}=\sum_\ell a_{j\alpha\ell}w^\ell$ and
$c_{j\alpha}\in\Q$. Let $\Omega_j=\{y:\eta_{y,0}\in W_j\}$, open and increasing, with union
$Y^\an$ since $\eta_{y,0}\in D$. On $\Omega_j$ put
$\phi_j=\max\{-j,\max_{\alpha,\ell\ge1}(\log|a_{j\alpha\ell}|-c_{j\alpha})/\ell\}$, which is convex PL
(terms vanishing at a point are dominated by $-j$ nearby). For $y\in\Omega_j$ and $r>0$,
Lemma~\ref{lem:gauss-point} gives $\eta_{y,r}\in W_j$ iff $\log r<-\phi_j(y)$, while
$\eta_{y,r}\in D$ iff $\log r<-u(y)$. From $W_j\subseteq W_{j+1}\subseteq D=\bigcup_jW_j$ we get
$\phi_j\ge\phi_{j+1}\ge u$ on $\Omega_j$, and for every rational $s>u(y)$ the point
$\eta_{y,e^{-s}}$ lies in some $W_j$, so $\phi_j(y)<s$; thus $\phi_j\downarrow u$ on $\Omega_j$. A
compact subset of $Y^\an$ lies in some $\Omega_{j_0}$, and Theorem~\ref{thm:import-local-to-global}
applies.
\end{proof}

\begin{lemma}[Principal-open hull inequality]\label{lem:principal-hull}
\lean{EnvelopeProperty.le_iSup_of_mem_hull_awayExtend, EnvelopeProperty.le_iSup_of_awayComap_mem_hull}
\leanok
\uses{lem:principal-extension, lem:hull-basic, lem:principal-open}
Let $0\ne f\in A$, $K\subseteq Y^\an$ compact and $x\in\widehat K\cap Y_f^\an$ (hull in $Y$). If
$v\in\Rcal(Y_f)$ and $v\le g+C$ for some $g\in\FS(Y)$, then $v(x)\le\sup_Kv$, where $v$ is extended by
$-\infty$ on $\{|f|=0\}$. (\cite{FHX} writes the conclusion as $v(x)\le\sup_{K\cap Y_f^\an}v$, which is the
same bound.)
\end{lemma}

\begin{proof}
Since $K$ is compact and $|f(x)|>0$, there is $M\ge0$ with $\log|f|\le M$ on $K$ and
$|\log|f(x)||\le M$. For $c\in\Q_{>0}$ let $v_c$ be the extension of $v+c\log|f|$ by $-\infty$ on
$\{|f|=0\}$; it is in $\Rcal(Y)$ by Lemma~\ref{lem:principal-extension}. On $K$ one has $v_c\le v+cM$
(both sides are $-\infty$ on $\{|f|=0\}$), so the hull inequality (Lemma~\ref{lem:hull-basic}) gives
\[
v(x)-cM\le v_c(x)\le\sup_Kv_c\le\sup_Kv+cM.
\]
Moreover $\sup_Kv\le\sup_Kg+C<\infty$, $g$ being continuous with values in $[-\infty,\infty)$. If
$\sup_Kv>-\infty$ let $c\downarrow0$; otherwise the inequality gives $v(x)=-\infty$.
\end{proof}

\subsection{A strictly psh exhaustion and gluing}\label{subsec:gluing}

\begin{lemma}[Strict exhaustion {\cite[Lemma~5.18]{FHX}}]\label{lem:strict-exhaustion}
\lean{EnvelopeProperty.exists_isStrictExhaustion, EnvelopeProperty.IsStrictExhaustion}
\leanok
\uses{lem:reg-closure-elementary, lem:reg-decreasing, lem:growth}
Let $Y$ be normal ($k$ countable). There is $\rho\in\Rcal(Y)$, $\rho\ge0$, such that for every
$x\in Y^\an$ and open $V\ni x$ there are $\rho'\in\Rcal(Y)$ and $\delta>0$ with $\rho'\le\rho$,
$\rho'(x)=\rho(x)$ and $\rho'\le\rho-\delta$ off $V$.
\end{lemma}

\begin{proof}
Let $\lambda(\tau)=\log(1+e^{-\tau})$ on $\mathbf T=\R\cup\{+\infty\}$, so that
$\lambda(-\log|f|)=\log(1+|f|)$. For $\tau^0\in\R$ let $p=e^{-\tau^0}/(1+e^{-\tau^0})\in(0,1)$ and
$\ell_{\tau^0}(\tau)=\lambda(\tau^0)-p(\tau-\tau^0)$; for $\tau^0=+\infty$ let $\ell_{\tau^0}=0$.
Then $\lambda-\ell_{\tau^0}\ge0$ with equality only at $\tau^0$, and it is bounded below by a positive
constant outside every neighbourhood of $\tau^0$: for $\tau^0=+\infty$ since $\lambda$ is positive and
decreasing; for $\tau^0\in\R$ since $\lambda-\ell_{\tau^0}$ is strictly convex with vanishing
derivative at $\tau^0$, tends to $+\infty$ at both ends ($\lambda(\tau)\ge\max\{0,-\tau\}$, $p<1$) and
equals $+\infty$ at $+\infty$.

From $\max_\alpha\binom m\alpha s^\alpha\le(1+s)^m\le(m+1)\max_\alpha\binom m\alpha s^\alpha$ for $s\ge0$, the
functions $P_m=\frac1m\max_{0\le\alpha\le m}(\log\binom m\alpha+\alpha\log|f|)$ satisfy
$P_m\le\log(1+|f|)\le P_m+\frac{\log(m+1)}m$; rounding constants gives regularizable
$Q_j$ (Lemma~\ref{lem:reg-closure-elementary}) with $\log(1+|f|)\le Q_j\le\log(1+|f|)+2^{-j-2}$, so
$Q_j+2^{-j}\downarrow\log(1+|f|)$ and $\log(1+|f|)\in\Rcal(Y)$ by Lemma~\ref{lem:reg-decreasing}. For
$\tau^0=-\log|f|(x)$ finite, $\ell_{\tau^0}(-\log|f|)=C+p\log|f|$ is the decreasing limit of
$\max\{C_m+a_m\log|f|,C_m+a'_m\log|f|,b/m-m\}\in\FS(Y)$ for rationals $a_m\uparrow p$,
$a'_m\downarrow p$, $C_m\downarrow C$, so it is regularizable.

Enumerate $A=\{f_1,f_2,\dots\}$ and put $h_j=\log(1+|f_j|)$. By Lemma~\ref{lem:growth} there are
integers $C_j\ge1$ with $0\le h_j\le C_jb$, $b=1+\max\{0,\log|z_l|\}$. Put $\gamma_j=2^{-j}/C_j$ and
$\rho=\sum_j\gamma_jh_j\le b$; the regularizable $R_N=\sum_{j\le N}\gamma_jh_j+2^{-N}b$ decrease to
$\rho$. Given $x$ and $V$, since the topology is the coarsest making all $|f|$ continuous there are
a finite $J$ and neighbourhoods $U_j$ of $\tau_j=-\log|f_j|(x)$ in $\mathbf T$ with
$\{y:-\log|f_j|(y)\in U_j\ \forall j\in J\}\subseteq V$. With $\delta_j>0$ a lower bound of
$\lambda-\ell_{\tau_j}$ off $U_j$, put
$\rho'=\rho-\sum_{j\in J}\gamma_j(\lambda-\ell_{\tau_j})(-\log|f_j|)$, $\delta=\min_j\gamma_j\delta_j$. Then
$\rho-\rho'\ge0$ vanishes at $x$ and is $\ge\delta$ off $V$, and
$R_N-\sum_{j\in J}\gamma_j(h_j-\ell_{\tau_j}(-\log|f_j|))$, $N\ge\max J$, are regularizable and decrease
to $\rho'$, which is finite at $x$; so $\rho'\in\Rcal(Y)$.
\end{proof}

\begin{lemma}[Hull criterion {\cite[Lemma~5.19]{FHX}}]\label{lem:hull-criterion}
\lean{EnvelopeProperty.IsRegularizable.of_hullIneq_add_smul}
\leanok
\uses{lem:hartogs, lem:strict-exhaustion, lem:hull-basic}
Let $u$ be upper semicontinuous of logarithmic growth, $u\not\equiv-\infty$,
$\widetilde u=u\circ p+\log|w|$ and $\widetilde\rho\ge0$ a finite function on $\widetilde Y^\an$ such that
every $c\widetilde\rho$, $c\in\Q_{>0}$, is upper semicontinuous (for instance $\widetilde\rho$ as in
Lemma~\ref{lem:strict-exhaustion} for $\widetilde Y$). If $\widetilde u+c\widetilde\rho$ satisfies the hull inequality for every rational
$c>0$, then $u\in\Rcal(Y)$.
\end{lemma}

\begin{proof}
For compact $K\subseteq\{\widetilde u+c\widetilde\rho<0\}$, $\sup_K(\widetilde u+c\widetilde\rho)<0$ by
upper semicontinuity, so $\widehat K\subseteq\{\widetilde u+c\widetilde\rho<0\}$. These Runge open sets
increase to $\{\widetilde u<0\}$ as $c\downarrow0$, since $\widetilde\rho$ is finite and nonnegative,
so $\{\widetilde u<0\}$ is Runge and Lemma~\ref{lem:hartogs} applies.
\end{proof}

\begin{proposition}[Gluing {\cite[Proposition~5.20]{FHX}}]\label{prop:gluing}
\lean{EnvelopeProperty.gluing_hullIneq, EnvelopeProperty.IsRegularizable.of_cover_principal}
\leanok
\uses{lem:strict-exhaustion, lem:rossi, lem:principal-hull, lem:reg-closure-elementary, lem:hull-criterion}
Let $Y$ be normal, $Z=V(I)\subsetneq Y$ closed, and $u$ upper semicontinuous of logarithmic growth
such that every $x\in Y^\an\setminus Z^\an$ lies in some $Y_g^\an$ with $u|_{Y_g^\an}\in\Rcal(Y_g)$.
Let $\rho$ be as in Lemma~\ref{lem:strict-exhaustion}. (1) For rational $c>0$, compact $K$ and
$x\in\widehat K\setminus Z^\an$, $(u+c\rho)(x)\le\sup_K(u+c\rho)$. (2) If $Z=\varnothing$ then
$u\in\Rcal(Y)$.
\end{proposition}

\begin{proof}
(1) Suppose $u_c=u+c\rho$ has $u_c(x)>\sup_Ku_c$. Let $v=\max_i\log|f_i|$ for generators $f_i$ of $I$
($v=0$ if $Z=\varnothing$), so $v\in\Rcal(Y)$ and $\{v=-\infty\}=Z^\an$. For small rational $\epsilon>0$
the strict inequality persists for $u_c+\epsilon v$, which attains its maximum over $\widehat K$ at
some $y\notin K\cup Z^\an$. Choose $Y_g\ni y$ with $u|_{Y_g^\an}$ regularizable, and a relatively
compact open $W\ni y$ with $\overline W\subseteq Y_g^\an$, $\overline W\cap K=\varnothing$; take
$\rho',\delta$ from Lemma~\ref{lem:strict-exhaustion}. Then $q=u+c\rho'+\epsilon v$ is regularizable
on $Y_g$ (Lemma~\ref{lem:reg-closure-elementary}), of logarithmic growth on $Y$, and on
$C=\partial W\cap\widehat K$ satisfies $\sup_Cq\le q(y)-c\delta$. But $y\in\widehat C$ by
Lemma~\ref{lem:rossi}, and Lemma~\ref{lem:principal-hull} gives $q(y)\le\sup_Cq$.
(2) Apply (1) with $Z=\varnothing$ to $\widetilde Y$, the cover $\widetilde Y_g$ and $u\circ p+\log|w|$,
regularizable on each $\widetilde Y_g$ by Lemma~\ref{lem:reg-closure-elementary}; then
Lemma~\ref{lem:hull-criterion} applies. $u\not\equiv-\infty$, being regularizable on some $Y_g$.
\end{proof}

\begin{remark}\label{rmk:gluing-hypotheses}
\uses{prop:gluing}
In \cite[Proposition~5.20]{FHX}, $u$ is weakly psh and $Y\setminus Z$ is covered by affine open
subsets on which $u$ is regularizable. Weakly psh functions are upper semicontinuous, which is the
only property used; and a cover by affine opens with $u$ regularizable yields one by principal
opens of $Y$ by Lemma~\ref{lem:reg-closure-elementary}(4). Covering every point of $Y^\an\setminus
Z^\an$ is equivalent to covering $Y\setminus Z$, every prime being the support of a point.
\end{remark}

\subsection{Weakly psh functions}\label{subsec:weakly-psh}

\begin{definition}[Weakly psh functions {\cite[Definition~5.6]{FHX}}]\label{def:weakly-psh}
\lean{EnvelopeProperty.IsWeaklyPSHOn, EnvelopeProperty.IsWeaklyPSHOn.upperSemicontinuousOn, EnvelopeProperty.isPshClass_isWeaklyPSHOn_univ}
\leanok
\uses{def:convex-pl}
Let $U\subseteq Y^\an$ be open. A function $u\colon U\to[-\infty,\infty)$ is \emph{weakly psh} if
there is a proper birational morphism $\mu\colon Y'\to Y$ such that $u\circ\mu^\an$ is locally psh on
$(\mu^\an)^{-1}(U)$. Weakly psh functions are upper semicontinuous, $\mu^\an$ being proper and
surjective.
\end{definition}

\begin{remark}\label{rmk:weakly-psh-blowup}
\uses{def:weakly-psh}
The formalization takes the witness $\mu$ to be the blowup of $Y=\Spec A$ along an ideal $(s)$, $s\subseteq A$
finite, covered by the affine charts $\Spec A[s/g]$, $g\in s\setminus\{0\}$; $u$ is weakly psh on $U$ when
$u\circ\mu^\an$ is locally psh on the preimage of $U$ in every chart. This is equivalent to the
intrinsic native proper-birational witness formulation above, as proved in
Theorem~\ref{paper:actual-proper-birational-weakly-psh-witness-iff}. The forward construction retains
the SAME original finite-set Rees scheme with actual proper and birational data. Conversely, actual
proper-birational Rees domination constructs an original finite family and a whole over-base map;
the original affine charts and exact scalar triangles then pull back the intrinsic locally psh
function to the unchanged chart definition. This equivalence is formalized for every domain $A$
and original subset $U$, without reducedness or projectivity of the proper-birational witness.
Classical psh and general-curve potential comparison remain separate analytic obligations.
Also formalized: the charts cover
the analytification and are proper over $Y^\an$ on the closed subsets $\{|h/g|\le1\}$ (surjectivity uses
a valuation ring dominating $A_{\mathfrak p}$), weakly psh functions are upper semicontinuous, and they
are stable under adding and taking maxima with regularizable functions.
\end{remark}

\subsection{Norms along generically \'etale projections}\label{subsec:norms}

\begin{lemma}[Finite flat norms]\label{lem:finite-flat-norm}
\lean{EnvelopeProperty.exists_multiset_algHom_algebraMap_norm_eq_prod, EnvelopeProperty.wgaussValuation_eq_prod_of_eval_eq_norm_finite, EnvelopeProperty.Berk.flatFibreMult, EnvelopeProperty.Berk.sum_flatFibreMult, EnvelopeProperty.Berk.sum_flatFibreMult_of_free, EnvelopeProperty.Berk.wgaussValuation_normForm_away, EnvelopeProperty.exists_free_localizationAway, EnvelopeProperty.flatNormFun, EnvelopeProperty.flatNormFun_eq_normFun, EnvelopeProperty.IsRegularizable.flatNormFun_of_free, EnvelopeProperty.IsRegularizable.flatNormFun}
\leanok
\uses{lem:reg-closure-elementary, lem:reg-decreasing, lem:reg-triv-finite, lem:finite-etale-norm, prop:gluing}
Let $\pi\colon Y\to S$ be a finite flat morphism of normal affine varieties of degree $d$. For
$s\in S^\an$ write $\mathcal H(s)\otimes_{k[S]}k[Y]=\prod_yA_y$, a product of local Artinian
$\mathcal H(s)$-algebras indexed by the points $y$ over $s$, and $m_y(\pi)=\dim_{\mathcal H(s)}A_y$,
so $\sum_{\pi(y)=s}m_y(\pi)=d$. Then $\Norm_\pi u(s)=\sum_{\pi(y)=s}m_y(\pi)u(y)$ maps $\Rcal(Y)$ to
$\Rcal(S)$. The base change statement of \cite[Lemma~5.13]{FHX} is Lemma~\ref{lem:finite-flat-norm-base-change}.
In Lean, $Y$ is only assumed integral (normality of $Y$ is not used), and $\sum_ym_y(\pi)=d$ is proved where
$k[Y]$ is free over $k[S]$; in general $\sum_ym_y(\pi)=\dim_{\mathcal H(s)}\mathcal H(s)\otimes_{k[S]}k[Y]$.
\end{lemma}

\begin{proof}
For $u=N^{-1}\max_j\{\log|b_j|+\lambda_j\}\in\FS(Y)$ ($b_j=1$ for the constant term), the
determinant of multiplication by $\sum_jT_jb_j$ on the projective $k[S][T]$-module $k[Y][T]$ is a
form $\sum_{|\alpha|=d}c_\alpha T^\alpha$ with $c_\alpha\in k[S]$, and
$N\Norm_\pi u=\max_{|\alpha|=d}\{\log|c_\alpha|+\sum_j\alpha_j\lambda_j\}$ on $S^\an$: after base change to
a complete algebraically closed field, the determinant on a local Artinian factor of length $\ell$ is
the $\ell$-th power of the residue of the linear form, and the weighted Gauss norm is
multiplicative. The right-hand side is regularizable by Lemma~\ref{lem:reg-closure-elementary}. For
$u\in\Rcal(Y)$ with approximants $u_i$, the $\Norm_\pi u_i$ decrease to $\Norm_\pi u$, which is finite
at $\vtriv$ (the only point over $\vtriv$ is $\vtriv$, with multiplicity $d$), so
Lemma~\ref{lem:reg-decreasing} applies.

In Lean the determinant on a projective non-free module is not available, so the norm forms are taken
over principal opens $S_h$ on which $k[Y]_h$ is free (they cover $S$, since a finite flat module over a
Noetherian ring is locally free), and the local regularizability statements are glued by
Proposition~\ref{prop:gluing} (this uses that $k$ is countable). The norm identity on a local Artinian factor
is proved after base change to an algebraically closed field: each local factor of $\Omega\otimes E$ has residue
field $\Omega$, and its norm is the $\dim$-th power of the residue, since the characteristic polynomial of a
nilpotent endomorphism is $X^{\dim}$.
\end{proof}

\begin{lemma}[Finite flat norms and base change {\cite[Lemma~5.13]{FHX}}]\label{lem:finite-flat-norm-base-change}
\lean{EnvelopeProperty.Berk.flatNormFun_scalarExtension}
\leanok
\uses{lem:finite-flat-norm,spectrum:paper-flat-norm-base-change}
In the setting of Lemma~\ref{lem:finite-flat-norm}, $(\Norm_\pi u)\circ r=\Norm_{\pi_F}(u\circ r)$ for complete
extensions $F/k$, where $r$ denotes the restriction maps from the analytifications over $F$.
\end{lemma}

\begin{proof}
The actual completed-residue tensor comparison identifies the new fibre algebra
with scalar extension of the original completed fibre. Its Artinian local-factor
dimensions are conserved after grouping maximal ideals by contraction. These
dimensions are precisely the multiplicities in both function norms. Regroup the
finite weighted sums to obtain the identity for every extended-real function.
The genuine tensor cancellation isomorphism supplies the literal
$F\otimes_k k[Y]$ presentation. The producer needs only module finiteness;
flatness and regularizability are not used in this pointwise calculation.

The formal comparison uses the $k$-trivial point spaces of the extended rings.
Restricting the base point to the prescribed norm of $F$ forces that same norm
on every point in its fibre, so this gives exactly the prescribed-$F$ statement.
\end{proof}

\begin{remark}\label{rmk:finite-flat-norm-etale}
\uses{lem:finite-flat-norm, prop:nonflat-norm}
Proposition~\ref{prop:nonflat-norm} uses Lemma~\ref{lem:finite-flat-norm} only over a principal open
$S_0$ over which $\pi$ is \'etale, and $S_0$ may be shrunk further so that $k[Y_0]$ is free over
$k[S_0]$. In that case $\mathcal H(s)\otimes_{k[S]}k[Y]$ is a finite product of finite separable field
extensions $L_y$ of $\mathcal H(s)$, $m_y(\pi)=[L_y:\mathcal H(s)]$, and the norm identity follows from
$\Norm_{L/\mathcal H(s)}=\prod_\sigma\sigma$ over the embeddings into an algebraic closure, the uniqueness of
the extension of the absolute value of the complete field $\mathcal H(s)$, and the multiplicativity of
weighted Gauss norms. The formalization proves this finite free \'etale form.
\end{remark}

\begin{lemma}[Finite free \'etale norms {\cite[Lemma~5.13]{FHX}}]\label{lem:finite-etale-norm}
\lean{EnvelopeProperty.Berk.fibreMult, EnvelopeProperty.Berk.sum_fibreMult, EnvelopeProperty.Berk.wgaussValuation_map_normForm, EnvelopeProperty.normFun, EnvelopeProperty.exists_normFun_eval_eq, EnvelopeProperty.IsRegularizable.normFun}
\leanok
\uses{lem:reg-closure-elementary, lem:reg-decreasing, lem:reg-triv-finite}
Let $k[S]\to k[Y]$ make $k[Y]$ a finite free \'etale $k[S]$-algebra of rank $d$, with $S$ normal and $Y$
integral. For $y\in Y^\an$ over $s$ put $m_y=\dim_{\mathcal H(s)}\mathcal H(s)\otimes_{k[S]}k[Y]/\mathfrak m_y$,
where $\mathfrak m_y$ is the kernel of $\mathcal H(s)\otimes_{k[S]}k[Y]\to\mathcal H(y)$. Then
$\sum_{\pi(y)=s}m_y=d$; for $b_1,\dots,b_N\in k[Y]$ and radii $r_j$, the weighted Gauss norm of the norm
form $\Norm(\sum_jT_jb_j)$ at $s$ is $\prod_{\pi(y)=s}(\max_j|b_j(y)|r_j)^{m_y}$; consequently
$\Norm_\pi$ maps finite maxima of terms $a\log|f|+b$ ($a>0$) to such maxima on $S^\an$, and maps
$\Rcal(Y)$ to $\Rcal(S)$.
\end{lemma}

\begin{proof}
\leanok
The fibre algebra is \'etale over the field $\mathcal H(s)$, hence a product of finite separable
extensions indexed by its maximal ideals, which correspond to the points over $s$; its dimension is
$d$. The norm form specializes to the norm of $\mathcal H(s)\otimes_{k[S]}k[Y]$, and the weighted
Gauss norm identity is the \'etale norm identity of Remark~\ref{rmk:finite-flat-norm-etale}. For an FS
function write it as $N^{-1}\max_j\{\log|b_j|+\lambda_j\}$ with $b_0=1$; taking logarithms gives
$N\Norm_\pi u=\max_\alpha\{\log|c_\alpha|+\langle\alpha,\lambda\rangle\}$, regularizable by
Lemma~\ref{lem:reg-closure-elementary}. For $u\in\Rcal(Y)$ the norms of the FS approximants decrease
to $\Norm_\pi u$, which is finite at $\vtriv$ (every point over $\vtriv$ is $\vtriv$), and
Lemma~\ref{lem:reg-decreasing} applies.
\end{proof}

\begin{lemma}[Subharmonic functions on curves {\cite[Lemma~5.21]{FHX}, \cite[Proposition~3.5]{GR25}}]\label{lem:curves}
\notready
\uses{def:weakly-psh}
Let $F$ be a complete non-Archimedean field. (1) If $q\colon C'\to C$ is a finite flat surjective
morphism of normal algebraic $F$-curves, $C$ integral, and $u$ is subharmonic on $C'^\an$, then
$\Norm_qu(y)=\sum_{q(x)=y}m_x(q)u(x)$ is subharmonic on $C^\an$. (2) Subharmonic functions on the
analytification of an algebraic curve that agree off finitely many rigid points are equal. (3) If
$F/k$ is algebraically closed and $u$ is weakly psh on $Y^\an$, then $u\circ r\circ p^\an$ is
subharmonic on $C^\an$ for every morphism $p\colon C\to Y_F$ from a smooth $F$-curve.
\end{lemma}

\begin{proof}
(1) Let $K=F(C)$, $E_i$ the function fields of the components of $C'$, $\nu_i$ the inseparable
degree of $E_i/K$, $L/K$ finite normal containing conjugates of all $E_i$, and $r_L\colon D\to C$
the normalization of $C$ in $L$. Each $K$-embedding $\sigma\colon E_i\to L$ gives
$p_{i,\sigma}\colon D\to C'_i$, and comparing characteristic polynomials of $b(T)=\sum_jT_jb_j$ for a
local basis $(b_j)$ shows $(\Norm_qu)\circ r_L^\an=\sum_i\nu_i\sum_\sigma u\circ p_{i,\sigma}^\an$. The right side
is subharmonic by the cone property and functoriality \cite[Proposition~3.5(4),(8)]{GR25}; since
$r_L^\an$ is closed and surjective, $\Norm_qu$ is upper semicontinuous and subharmonic by finite
surjective descent \cite[Proposition~3.5(8)]{GR25}. (2) At a rigid point $z$, upper semicontinuity
gives $h(z)\ge\limsup_{x\to z,x\ne z}h(x)$, and strict inequality would be a strict local maximum,
contradicting the maximum principle \cite[Proposition~3.5(6)]{GR25}. (3) For a witness $\mu$,
$u\circ r\circ\mu_F^\an$ is locally psh, hence classically psh \cite[Proposition~4.2]{GR25}. A component of
$C\times_{Y_F}Y'_F$ dominating $C$ gives, after normalization, a smooth curve $C''$ finite and
surjective over $C$ with a lift $C''\to Y'_F$; the pullback to $C''^\an$ is subharmonic and finite
surjective descent applies.
\end{proof}

\begin{lemma}[Extension step of Proposition~\ref{prop:nonflat-norm} {\cite[Proposition~5.22]{FHX}}]\label{lem:nonflat-extension}
\lean{EnvelopeProperty.IsRegularizable.exists_extend, EnvelopeProperty.IsRegularizable.exists_extend_of_hasEnvelopeProperty, EnvelopeProperty.IsRegularizable.exists_extend_of_charP, EnvelopeProperty.IsRegularizable.exists_extend_of_isAffineSpacePrincipalOpen, EnvelopeProperty.IsRegularizable.exists_extend_of_isAffineSpacePrincipalOpen_of_charP, EnvelopeProperty.ProjBerk.exists_extend_mulChart, EnvelopeProperty.ProjBerk.exists_eval_le_chartLog}
\leanok
\uses{lem:principal-extension, lem:growth, lem:bj-chart, lem:bj-compactification-regularize, lem:bj-compactification-extend, def:pn-ep, thm:bj-psh2, thm:bj-psh-basic, thm:pn-ep-charp}
Let $k$ be countable, $n\ge1$, $0\ne h_S\in k[y_1,\dots,y_n]$, $S=\Spec k[y]_{h_S}$, $0\ne h\in k[S]$ and $S_0=S_h$. Assume that
every rational $\theta>0$ has the envelope property on $\PP^n$ (Theorem~\ref{thm:pn-ep-charp} in characteristic $p$) and
the imported Theorem~\ref{thm:import-growth-control}. Let $v_0\in\Rcal(S_0)$ with $v_0\le g+B$ on $S_0^\an$ for some
$g\in\FS(S)$ and $B\in\R$. Then there is $v\in\Rcal(S)$ with $v=v_0$ on $S_0^\an$.
\end{lemma}

\begin{proof}
(1) Let $\tilde h_S\in k[x_0,\dots,x_n]$ be a form of degree $e$ with $\tilde h_S(1,y)=h_S$ (e.g.\ the homogenization) and $\sigma=x_0\tilde h_S$, so that
$S\cong P_\sigma$ and $S^\an\cong U_\sigma$ (Lemma~\ref{lem:bj-chart}(1), the chart $x_0\cdot\tilde h_S$). (2) Every element of
$A_\sigma=k[S]$ is $\tau/\sigma^N$ with $|\tau|\le1$, so $\log|\tau/\sigma^N|\le-N\log|\sigma|$; hence every $g\in\FS(S)$ satisfies
$g\le-\kappa_g\log|\sigma|+C_g$ for a rational $\kappa_g>0$, and $\log|h|\le-N_h\log|\sigma|$. Note $-\log|\sigma|\ge0$. (3) Let
$c_i=1/(i+1)$. By Lemma~\ref{lem:principal-extension} the function $w_i$ equal to $v_0+c_i\log|h|$ on $S_0^\an$ and to $-\infty$
on $S^\an\setminus S_0^\an$ lies in $\Rcal(S)$, and $w_i\le-\kappa\log|\sigma|+B_0$ with $\kappa=\kappa_g+N_h$ and
$B_0=B+C_g$ by (2), since $c_i\log|h|\le c_i(-N_h\log|\sigma|)\le-N_h\log|\sigma|$. (4) Fix a rational $a>\kappa$. By
Lemma~\ref{lem:bj-compactification-extend}, $\phi_i\in\PSH(a(e+1))$, where $\phi_i=w_i+a\log|\sigma|$ on $U_\sigma$ and
$\phi_i=-\infty$ elsewhere; and $\phi_i\le B_0+(a-\kappa)\log|\sigma|\le B_0$. (5) By the envelope property
$\psi_j=(\sup_{i\ge j}\phi_i)^\star\in\PSH(a(e+1))$; the $\psi_j$ decrease, and their limit $\psi$ is in
$\PSH(a(e+1))$ (Theorem~\ref{thm:bj-psh-basic}, Theorem~\ref{thm:bj-psh2}) provided $\psi\not\equiv-\infty$, which follows
from (6) at $\vtriv\in S_0^\an$. (6) On the open set $S_0^\an$, for $i\ge j$, $c_i\log|h|\le c_j\max(\log|h|,0)$, so
$\phi_i\le v_0+c_j\max(\log|h|,0)+a\log|\sigma|$, an upper semicontinuous function on $S_0^\an$; hence $\psi_j$ is bounded by it on
$S_0^\an$, and $\psi\le v_0+a\log|\sigma|$ there. Conversely $\psi_j\ge\phi_i$ for $i\ge j$, and $\phi_i\to v_0+a\log|\sigma|$ on
$S_0^\an$, where $\log|h|$ is finite. So $\psi=v_0+a\log|\sigma|$ on $S_0^\an$. (7) By Lemma~\ref{lem:bj-compactification-regularize},
$v=(\psi-a\log|\sigma|)|_{S^\an}\in\Rcal(S)$, and $v=v_0$ on $S_0^\an$.
\end{proof}

\begin{proposition}[Nonflat norms {\cite[Proposition~5.22]{FHX}}]\label{prop:nonflat-norm}
\lean{EnvelopeProperty.NonflatNorms}
\notready
\uses{lem:finite-etale-norm, lem:finite-flat-norm, lem:finite-flat-norm-base-change, lem:curves, prop:gluing, lem:nonflat-extension, lem:growth, lem:reg-triv-finite, thm:pn-ep, spectrum:paper-rigid-complete-extension, algebra:affine-line-avoids-polynomial, scheme:affine-line-closed-immersion}
Let $Y$ be a normal affine variety of dimension $n\ge1$, $S\subseteq\A^n$ a principal open subset,
$\pi\colon Y\to S$ finite surjective generically \'etale of degree $d$, and $Z\subsetneq Y$ closed.
There are integers $m_y(\pi)\ge1$, $y\in Y^\an$, with $\sum_{\pi(y)=s}m_y(\pi)=d$ for all $s$, such
that $\Norm_\pi u\in\Rcal(S)$ for every weakly psh $u$ of logarithmic growth that is regularizable on
an affine open cover of $Y\setminus Z$.
\end{proposition}

\begin{proof}
Choose $0\ne h\in k[S]$ with $S_0=S_h\subseteq S\setminus\pi(Z)$ and $\pi$ \'etale over $S_0$; over
$S_0$ use the weights of Lemma~\ref{lem:finite-flat-norm}. For $s\notin S_0^\an$ let $F$ be the
completion of an algebraic closure of $\mathcal H(s)$ and $s_F$ the induced rigid point; choose a line
$\Lambda\ni s_F$ not contained in the closure of $(S\setminus S_0)_F$, put $C_s=\Lambda\cap S_F$ and let
$C'_s$ be the normalization of the reduction of $Y_F\times_{S_F}C_s$; then $q_s\colon C'_s\to C_s$ is
finite flat of degree $d$ between smooth curves, an isomorphism onto the fibre product over
$(S_0)_F$. Define $m_y(\pi)$ as the sum of $m_x(q_s)$ over $x\in q_s^{-1}(s_F)$ mapping to $y$.

Regularity: $u|_{Y_0}\in\Rcal(Y_0)$ by Proposition~\ref{prop:gluing}(2), so
$v_0=\Norm_{\pi|_{Y_0}}u\in\Rcal(S_0)$ by Lemma~\ref{lem:finite-flat-norm}, and
$v_0\le dB'g+dB''$ with $g\in\FS(S)$ by Lemma~\ref{lem:growth}. Extension (Lemma~\ref{lem:nonflat-extension}): with $\PP^n\supseteq\A^n$,
$\sigma\in H^0(\PP^n,\mathcal O(e+1))$ with $P_\sigma=S$, and rationals $c_i\downarrow0$, the extensions
$w_i$ of $v_0+c_i\log|h|$ lie in $\Rcal(S)$ (Lemma~\ref{lem:principal-extension}) with
$w_i\le-\kappa\log|\sigma|+B_0$; their compactified versions $\phi_i\in\mathrm{PSH}(\PP^n,aH)$ satisfy
$\phi_i\le B_0$, and by the envelope property of $\PP^n$ (FHX Theorem~3.1) the functions
$\psi_j=\sup^*_{i\ge j}\phi_i$ are $aH$-psh and decrease to $\psi$ with $\psi=v_0+a\log|\sigma|$ on
$S_0^\an$; $v=(\psi-a\log|\sigma|)|_{S^\an}\in\Rcal(S)$ and $v=v_0$ on $S_0^\an$. Identification:
$v\circ r$ and $\Norm_{q_s}$ of the pullback of $u\circ r$ are subharmonic on $C_s^\an$
(Lemma~\ref{lem:curves}), agree on $(C_s\cap(S_0)_F)^\an$ by
Lemma~\ref{lem:finite-flat-norm-base-change}, hence everywhere by Lemma~\ref{lem:curves}(2); evaluating at $s_F$
gives $v(s)=\Norm_\pi u(s)$.
\end{proof}

\subsection{Extension across a closed subset}\label{subsec:extension}

\begin{lemma}[Projections {\cite[Lemma~5.23]{FHX}}]\label{lem:projection-existence}
\lean{EnvelopeProperty.ProjectionExistence, EnvelopeProperty.exists_projection_of_isAlgClosed, EnvelopeProperty.projectionExistence_of_isAlgClosed}
\leanok
Let $Y$ be an affine variety of dimension $n\ge1$, $Z\subsetneq Y$ closed and $z\in Z$ a closed point
at which $Z$ is smooth. There are an affine open $U\ni z$, a principal open $S\subseteq\A^n$ and a
finite surjective generically \'etale $\pi\colon U\to S$ restricting to a closed immersion on
$Z\cap U$.
\end{lemma}

\begin{proof}
Compose general linear projections of the closure $X\subseteq\PP^N$ of $Y$ from closed points, so that
the composite $r\colon X\to\PP^n$ is finite and restricts to a closed immersion of
$B=\overline Z\cup(X\setminus Y)$ over a neighbourhood $N_z$ of $r(z)$ (at each step
\cite[Tag~0B1Q]{Stacks} applies to the component of the image of $B$ through the image of $z$, and
the centre avoids the cones joining that point to the other components). Fix a smooth closed point
$w\in X$ and choose each centre so that the projection is a local isomorphism at the image of $w$
(unramified at the last step); then $r$ is unramified at $w$, hence generically \'etale. Take a
principal open $S\ni r(z)$ of a chart $\A^n$ with $S\subseteq N_z\setminus r(X\setminus Y)$,
$U=r^{-1}(S)$ and $\pi=r|_U$.
\end{proof}

\begin{proposition}[Projection step {\cite[Proposition~5.24]{FHX}}]\label{prop:projection-step}
\lean{EnvelopeProperty.projection_hullIneq, EnvelopeProperty.IsRegularizable.of_projection, EnvelopeProperty.IsRegularizable.of_projection_weaklyPSH}
\uses{prop:nonflat-norm, lem:strict-exhaustion, prop:gluing, lem:rossi, lem:hull-criterion, lem:reg-closure-elementary, cor:analytic-closed-restriction-injective}
Let $Y,S,\pi,d$ be as in Proposition~\ref{prop:nonflat-norm} and $Z\subsetneq Y$ closed such that
$\pi|_Z$ is a closed immersion and $Z'=\pi^{-1}(\pi(Z))\setminus Z$ is closed. If $u$ is weakly psh
of logarithmic growth and regularizable on an affine open cover of $Y\setminus Z$, then
$u\in\Rcal(Y)$.
\end{proposition}

\begin{proof}
By Lemma~\ref{lem:hull-criterion} applied on $\widetilde Y$ (with $Z\times\A^1$, $\pi\times\mathrm{id}$,
$u\circ p+\log|w|$, which satisfy the same hypotheses), it suffices to show that $u_c=u+c\rho$
satisfies the hull inequality for every rational $c>0$, where $\rho$ is a strict exhaustion. Suppose
$u_c(x_0)>\sup_Lu_c$ with $x_0\in\widehat L$. Let $x$ be a maximum point of $u_c$ on $\widehat L$;
normalize $u_c(x)=0=\max_{\widehat L}u_c$ and choose rational $\delta>0$ with $\sup_Lu_c\le-\delta$, so
$x\notin L$ and $x\in Z^\an$ by Proposition~\ref{prop:gluing}(1). Choose $f_1+f_2=1$ with $f_1=0$ on
$Z'$ and $f_2=0$ on $Z$; then $O_1=\{|f_2|<1\}\supseteq Z^\an$ and $O_2=\{|f_1|<1\}\supseteq Z'^\an$ are
disjoint, $|f_1|=1$ on $O_1$, $|f_2|=1$ on $O_2$. The restriction of $\pi^{\an}$ to $Z^{\an}$
is injective by Corollary~\ref{cor:analytic-closed-restriction-injective}.
The fibre $\pi^{-1}(\pi(x))$ is $\{x\}\cup\Gamma$ with
$\Gamma\subseteq Z'^\an$ compact. Choose: a relatively compact open $A\supseteq\Gamma$ with
$\overline A\subseteq\{\log|f_1|<-1\}$ and rational $M\ge0$ with $u_c\le M$ on $\overline A$; an open
$V_x\ni x$ inside $\{\log|f_2|<-1\}$ with $u_c<\epsilon=\delta/(2d)$; $V=S^\an\setminus
\pi(Y^\an\setminus(V_x\cup A))$, so $\pi^{-1}(V)\subseteq V_x\cup A$; a relatively compact open
$B\ni x$ with $\overline B\subseteq V_x\cap\pi^{-1}(V)$, $\overline B\cap L=\varnothing$; and
$\rho',\delta'$ from Lemma~\ref{lem:strict-exhaustion} for $x$ and $B$. Put $u'=u+c\rho'$,
$\eta=\max\{\delta\log|f_2|,-\delta\}$ and $u^\sharp=\max\{u'+(M+1)\log|f_1|,\eta\}$, which is weakly psh,
of logarithmic growth and regularizable where $u$ is. On $V_x$, $u^\sharp=\max\{u',-\delta\}<\epsilon$;
on $A$, $u^\sharp=0$. So $h=\Norm_\pi u^\sharp\in\Rcal(S)$ (Proposition~\ref{prop:nonflat-norm}) and
$h(\pi(x))=0$. By Lemma~\ref{lem:rossi}, $x$ is in the hull of $K=\partial B\cap\widehat L$, so $\pi(x)$ is
in the hull of $\pi(K)$ and $0\le\sup_Kh\circ\pi$. For $y\in K$: if $y\notin Z^\an$ then
$u^\sharp(y)=-\delta$ and $h(\pi(y))\le-\delta m_y+(d-m_y)\epsilon<-\delta/2$; if $y\in Z^\an$ the other
points of the fibre lie in $A$ and $h(\pi(y))=m_y\max\{u'(y),-\delta\}\le m_y\max\{-c\delta',-\delta\}<0$.
Upper semicontinuity of $h\circ\pi$ gives a contradiction.
\end{proof}

\begin{remark}\label{rmk:projection-step-conditional}
\uses{prop:projection-step}
The proof of Proposition~\ref{prop:projection-step} uses weak plurisubharmonicity of $u$ only through
upper semicontinuity, the stability of the class under adding regularizable functions and nonnegative
rational multiples of $\log|f|$, under maxima with regularizable functions and real translates, and
through Proposition~\ref{prop:nonflat-norm} applied to $u^\sharp$. The formalization isolates exactly
these properties: it proves the barrier argument for any class of functions with these closure
properties, given weights as in Proposition~\ref{prop:nonflat-norm} for $\pi$ and for
$\pi\times\mathrm{id}_{\A^1}$. In Lean, the class is \texttt{IsPshClass} (upper semicontinuity and
stability under adding, and taking maxima with, regularizable functions; multiples of $\log|f_1|$ and
constants are regularizable), the output of Proposition~\ref{prop:nonflat-norm} is the structure
\texttt{NormWeights}, $Z=V(S)$, ``$\pi|_Z$ is a closed immersion'' is surjectivity of
$k[S]\to k[Y]/(S)$, and ``$Z'$ closed'' is encoded by the partition $f_1+f_2=1$. The barrier argument
is \texttt{projection\_hullIneq}; \texttt{IsRegularizable.of\_projection} concludes through
Lemma~\ref{lem:hull-criterion}, taking the class and the weights on $\widetilde Y\to S\times\A^1$ as
hypotheses and deriving the rest. For weakly psh functions (Definition~\ref{def:weakly-psh}, blowup witnesses),
\texttt{IsRegularizable.of\_projection\_weaklyPSH} derives all data on $\widetilde Y$ (principal open of
$\A^{n+1}$, normality of $k[Y][w]$, finiteness and generic \'etaleness of $\pi\times\mathrm{id}$, weak
plurisubharmonicity of $u\circ p+\log|w|$) and uses Proposition~\ref{prop:nonflat-norm} through the
explicit hypothesis class \texttt{NonflatNorms} (the author's decision: imported for now, to be
proved). So Proposition~\ref{prop:projection-step} depends only on \texttt{NonflatNorms} and
\texttt{WuRegularization}.
\end{remark}

\begin{theorem}[Local regularization {\cite[Theorem~5.25]{FHX}}]\label{thm:normal-local-regularization}
\lean{EnvelopeProperty.IsRegularizable.of_isWeaklyPSHOn_of_localPositivity, EnvelopeProperty.IsRegularizable.of_isWeaklyPSHOn_of_isAlgClosed, EnvelopeProperty.IsRegularizable.of_isWeaklyPSHOn, EnvelopeProperty.IsWeaklyPSHOn.comp_awayComap}
\uses{prop:gluing, lem:projection-existence, prop:projection-step, lem:reg-closure-elementary}
Let $Y$ be a normal affine variety over the countable field $k$, $Z\subsetneq Y$ closed and $u$
weakly psh of logarithmic growth, regularizable on an affine open cover of $Y\setminus Z$. Then
$u\in\Rcal(Y)$.
\end{theorem}

\begin{proof}
Let $\Sigma$ be the union of the affine opens on which $u$ is regularizable and $W=Y\setminus\Sigma\subseteq Z$;
by Proposition~\ref{prop:gluing}(2) it suffices to show $W=\varnothing$. Otherwise $n=\dim Y\ge1$,
$W$ (reduced, $k$ perfect) is smooth at a closed point $z$, and there is a principal open
$Y_1\ni z$ with $W\cap Y_1$ smooth and irreducible. Lemma~\ref{lem:projection-existence} gives
$\pi\colon U\to S$ restricting to a closed immersion on $W_U=W\cap U$. With
$Z'=\pi^{-1}(\pi(W_U))\setminus W_U$ and $E=\overline{Z'}\cap W_U$ a proper closed subset of $W_U$,
choose $g\in k[S]$ vanishing on $\pi(E)$ but not on $\pi(W_U)$, and put $S_1=S_g$, $U_1=\pi^{-1}(S_1)$.
Then $U_1\to S_1$ and $W_1=W\cap U_1\ne\varnothing$ satisfy the hypotheses of
Proposition~\ref{prop:projection-step}, so $u|_{U_1}$ is regularizable and $U_1\subseteq\Sigma$, a
contradiction.
\end{proof}

\begin{remark}\label{rmk:local-regularization-conditional}
\uses{thm:normal-local-regularization}
The formalized Theorem~\ref{thm:normal-local-regularization} (weakly psh functions with blowup
witnesses; ``regularizable on an affine open cover of $Y\setminus Z$'' in the equivalent principal-open
form) assumes \texttt{WuRegularization} (Remark~\ref{rmk:imports}),
\texttt{NonflatNorms} (Proposition~\ref{prop:nonflat-norm}) and \texttt{ProjectionExistence}, which is exactly Lemma~\ref{lem:projection-existence} (with a general
affine open $U$); the latter is proved for every algebraically closed field
(\texttt{projectionExistence\_of\_isAlgClosed}, by an affine route: Noether coordinates with a Nagata
substitution, a cotangent bound from smoothness of $Z$ at $z$, twisted coordinates $t_i^{f_i}-a_i$, generic
\'etaleness by perturbing the $a_i$ inside $I$, and Nakayama), so
\texttt{IsRegularizable.of\_isWeaklyPSHOn\_of\_isAlgClosed} assumes only \texttt{WuRegularization} and
\texttt{NonflatNorms}. Generic smoothness, the shrinking of $U$ to a principal open, the choice of $g$ and the
partition of unity separating $W_1$ from $Z'_1$ are proved
(\texttt{ProjectionExistence.toProjectionPackage}); the choice of $g$ uses incomparability for the finite
$\pi$ (every minimal prime of $IA_a$ is minimal over $\pi^{-1}(IA_a)A_a$) instead of the dimension argument,
and irreducibility of $W\cap Y_1$ is not needed. The two remaining hypotheses are to be proved.
\end{remark}

\begin{corollary}[Psh descent along modifications]\label{cor:normal-weak-global}
\lean{EnvelopeProperty.ProjSubvariety.isPSH_of_isPSH_comp_proj, EnvelopeProperty.ProjSubvariety.isPSH_of_isPSH_comp_proj_of_convexPLApproximation}
\leanok
\uses{thm:normal-local-regularization, thm:pg-psh-descent, rmk:pg-psh-descent-source}
Let $X$ be a normal projective variety over the countable field $k$, $L$ an ample $\Q$-line bundle,
and $\mu\colon X'\to X$ birational from an integral projective $X'$. If $\varphi\circ\mu^\an\in
\mathrm{PSH}(X',\mu^*L)$ then $\varphi\in\mathrm{PSH}(X,L)$.
\end{corollary}

\begin{proof}
Replacing $X'$ by its normalization $\nu\colon X'^\nu\to X'$ preserves the hypothesis, since psh functions pull back
(\cite[Proposition~2.1(4)]{FHX}: $\varphi\circ(\mu\nu)^\an\in\mathrm{PSH}(X'^\nu,(\mu\nu)^*L)$), and $\mu\circ\nu$ is again
birational; so we may assume $X'$ normal. Then, as in \cite[Corollary~5.26]{FHX}: reduce to $L$ very ample with sections $z_i$; on $U_i=X_{z_i}$ the
function $u_i=\varphi-\log|z_i|$ is weakly psh with witness $\mu$, of logarithmic growth, and
regularizable on affine opens covering $U_i\setminus Z$ ($Z$ the exceptional locus of $\mu$), by the
compactification lemma on $X'$; Theorem~\ref{thm:normal-local-regularization} gives $u_i\in\Rcal(U_i)$,
the compactification lemma gives $\psi_{i,a}\in\mathrm{PSH}(X,aL)$ for rational $a>1$, $\max_i\psi_{i,a}=\varphi$,
and PSH2 as $a\downarrow1$ concludes. In the model of Section~\ref{sec:projective-general} this is
Theorem~\ref{thm:pg-psh-descent}; see Remark~\ref{rmk:pg-psh-descent-source}.
\end{proof}
